\documentclass[12pt]{article}
\usepackage[UTF8]{ctex}
\usepackage[a4paper,margin=2.2cm]{geometry}
\usepackage{setspace}
\setstretch{1.25}
\title{《AI算力中心任务调度问题》讲题讲稿（约20分钟）}
\author{配套 slides.tex}
\date{}
\begin{document}
\maketitle

\section*{使用说明}
本讲稿按页对应 PPT（共 15 页内容页，不含封面目录），总时长约 20 分钟。
每页给出：\textbf{建议时长}、\textbf{讲述要点}、\textbf{可直接念的串词}。

\section*{第1页：背景与建模对象（1.5分钟）}
\textbf{串词：}
“这道题的业务背景很完整：校内 AI 平台给同学提供免费算力，任务会不断到达。系统不是只有一个瓶颈，而是三个并行约束：最多并发进程数 n、GPU 总量 K、内存总量 M。每个任务都有自己的资源需求和重要性。我们要做的是一个在线调度器——事件来了立刻做决策，不能等待未来信息。”

\section*{第2页：任务模型与目标函数（2分钟）}
\textbf{串词：}
“题目里最关键的不确定性是 checkpoint 总数。每个任务的 k\_i、m\_i、t\_i、W\_i 都可见，但 c\_i 在任务结束前未知。并且题目约定总时长就是 c\_i 乘 t\_i。目标函数是 $\sum W_i\cdot T_i/(t_i c_i)$，也就是在单位工作量意义下压缩高重要性任务的周转时间。所以这题天然不可能做离线精确最优，必须做在线估计和启发式决策。”

\section*{第3页：checkpoint 与抢占坑点（2分钟）}
\textbf{串词：}
“checkpoint 语义是判题陷阱。任务开始时刻就算第一个 checkpoint，结束时刻不算 checkpoint；抢占本身不扣时间，但会清空最近一段未存档进度。也就是说，任务刚过 checkpoint 时被抢占损失较小，快到下个 checkpoint 时被抢占损失很大。我们后面的抢占策略必须显式比较收益和回滚损失，而不是只看任务权重。”

\section*{第4页：交互机制、合法性与评分（2.5分钟）}
\textbf{串词：}
“实现层面是交互式 API：Init 做初始化，所有决策都在 Event 里触发。事件有四类：checkpoint、完成、新任务、预约调度。注意同一时刻可能触发多个 Event 调用，你不能假设一次性拿到完整批次。还有两个必须强调的合法性：Schedule 时任务不能已经在别的进程跑、且资源必须够；Appoint 的时间不能早于 now\_time。违反这些直接 0 分。评分则按和标程目标值比值给分，且每个测试点只有 8 秒时限。”

\section*{第5页：总体策略（1.5分钟）}
\textbf{串词：}
“题解代码可以概括为三步：先给任务打分，再在资源约束下选集合 S，最后对集合 S 做列表调度。如果某任务本来就在跑而且还在 S 里，尽量不动；只对需要新启动的任务考虑空闲进程或抢占。这个结构简单但实用，复杂度可控，适合 1 万任务级别在线调度。”

\section*{第6页：数据结构（1.5分钟）}
\textbf{串词：}
“Task 结构里，静态信息是 k,m,t,W，动态信息包括状态、已过 checkpoint 数、等待时间累计、当前绑定进程。Process 结构只记当前任务和当前段开始时间。全局有 usedK/usedM 做资源账本，还有 mean\_segments 作为未知 checkpoint 总数的先验估计。这个状态设计确保每次事件更新都能 O(1) 地改关键字段。”

\section*{第7页：打分公式（2分钟）}
\textbf{串词：}
“打分先估计任务总段数，$\hat c_i=passed+mean\_segments$，再算单位价值 $\phi_i=W_i/(\hat c_i t_i)$。然后除以资源占用强度和时长惩罚，得到最终分数 sc。直观上，高权重、预计更短、占资源更省、单段更短的任务优先。这里还有一个亮点：资源强度里有 $\alpha$，会根据当前 GPU/内存紧张程度自适应偏置。”

\section*{第8页：选集合 S（1.5分钟）}
\textbf{串词：}
“如果直接全排序再精确求最优子集，代价太高。题解用 nth\_element 先截出高分候选，再按分数降序做三约束贪心装箱：进程、GPU、内存都要同时满足。这本质上是多约束背包的近似，虽非最优，但在线环境里性价比很高。”

\section*{第9页：抢占判定（2分钟）}
\textbf{串词：}
“抢占判断不是拍脑袋，而是算不等式：新任务收益是否大于旧任务已投入损失乘以阈值。旧任务有两个量：已经跑了多久 s\_j（会丢），离下个 checkpoint 还有多久 r\_j（放弃后续机会成本）。代码同时构造一个 evict\_cost，在可抢占进程里选代价最小者。这样能避免过度抖动。”

\section*{第10页：事件处理（1.5分钟）}
\textbf{串词：}
“checkpoint 事件就给任务过段数加一，并重置段起点；结束事件要释放资源、更新 mean\_segments 并删任务；到达事件创建新任务入池；预约事件只触发调度。最后根据事件计数决定是否这次真正重算调度，兼顾响应性和耗时。”

\section*{第11页：合法性保证（1.5分钟）}
\textbf{串词：}
“这题判 0 分的高风险点是非法调用。题解里先做 task\_process\_check，再执行调度，确保资源和进程范围合法；任务状态机保证不会出现同一 task 同时在多个进程运行；Appoint 也保证时间不回退。建议比赛时把这些 assert 风格检查都保留到最后再按性能考虑裁剪。”

\section*{第12页：时限把控（2分钟）}
\textbf{串词：}
“最大 1 万任务、1000 进程、8 秒时限，必须做工程优化。题解用了哈希 + 向量、局部排序、抽样重调度。尤其是事件很密时，不是每次都全量重排，而是前期全调度，后期每 4 次调 1 次，但到达/结束强制调度。这个策略在真实日志类数据上通常更稳。”

\section*{第13页：样例推演（1.5分钟）}
\textbf{串词：}
“样例里两个任务资源冲突，不能并跑。策略就是 task1 从 0 跑到 6，再切 task2。得到目标值 19.1667。这个例子强调一点：有时‘不抢占’反而更优，因为抢占会丢失段内进度，而收益不一定覆盖损失。”

\section*{第14页：参数与优化方向（1.5分钟）}
\textbf{串词：}
“题解表现很依赖参数：短段偏好、抢占阈值、驱逐损失权重、段数先验更新速度。建议把测试集按‘真实日志风格’和‘理想分布风格’分组，分别看抢占频率、平均回退损失和长尾等待，再做参数微调，通常比盲调目标值更有效。”

\section*{第15页：总结与答疑（1分钟）}
\textbf{串词：}
“最后总结：这题难点不是某个复杂算法，而是在线信息不全下的稳定决策。题解通过‘估计段数 + 多资源打分 + 成本化抢占’把问题拆成可实现、可调参、可过时限的工程方案。我的建议是先保证合法性和稳定性，再逐步调参数冲分。谢谢大家，欢迎提问。”

\end{document}
